%============================ RESUMEN =======================================
\chapter*{Resumen ejecutivo}
\addcontentsline{toc}{chapter}{Resumen ejecutivo}
\markboth{RESUMEN EJECUTIVO}{}

El presente proyecto documenta el diseño, la implementación y la validación de \textbf{EILOR}, un sistema integral de \emph{monitoreo del espectro radioeléctrico de 2.4~GHz} y \emph{gestión de acceso de invitados} orientado a pequeños negocios y a laboratorios de aprendizaje en ciberseguridad. El sistema combina un dispositivo de hardware de bajo costo basado en el microcontrolador \textbf{ESP32}, que integra ambas radios y escanea redes Wi\nobreakdash-Fi (canales 1 al 13) y balizas Bluetooth de Baja Energía (BLE) en modos dedicados conmutables, con un \textbf{servidor de tiempo real} construido sobre Node.js, Express y WebSocket que agrega, analiza, persiste y visualiza los datos capturados en un tablero de mando (\emph{dashboard}) web.

EILOR incorpora cuatro modos operativos ---análisis de espectro Wi\nobreakdash-Fi, radar de proximidad BLE por RSSI, punto de acceso (SoftAP) con \emph{portal cautivo} de registro de invitados y modo de telemetría/reposo--- gobernados de forma remota mediante comandos WebSocket. En la capa de análisis, el servidor calcula un \emph{puntaje de interferencia} por canal que pondera el solapamiento de canales adyacentes y la intensidad de señal, recomienda el canal más limpio del conjunto no solapado 1/6/11 y detecta puntos de acceso sospechosos de tipo \emph{gemelo malicioso} (\emph{evil twin}). Adicionalmente, se implementó una capa de seguridad con autenticación por token de sesión para las acciones de control, un token opcional de dispositivo para restringir la ingesta de datos y un mecanismo de identificación de fabricante mediante prefijos OUI, junto con la detección de direcciones MAC aleatorizadas por privacidad.

La validación funcional confirmó la operación de extremo a extremo del sistema: ingesta de escaneos, análisis de canales, persistencia con recuperación tras reinicio, control de acceso (respuestas 401/200 según credenciales) y registro de invitados desde el portal cautivo con \emph{buffer} de reenvío ante caídas de enlace. El documento presenta, además, el análisis de desempeño, la estimación de costos con curva~S, la evaluación de impacto ambiental, la propuesta de escalabilidad y los manuales de usuario y técnico, cumpliendo la totalidad de los entregables definidos en la macro actividad.

\vspace{0.6cm}
\noindent\textbf{Palabras clave:} espectro 2.4~GHz, ESP32, IEEE 802.11, Bluetooth Low Energy, WebSocket, portal cautivo, ciberseguridad, gemelo malicioso, OUI, tiempo real.

\vspace{0.4cm}
\hrule
\vspace{0.4cm}

\noindent\textbf{\emph{Abstract.}} This project documents the design, implementation and validation of \textbf{EILOR}, an integrated system for \emph{2.4~GHz radio-frequency spectrum monitoring} and \emph{guest-access management} aimed at small businesses and cybersecurity learning labs. It couples a low-cost \textbf{ESP32}-based device ---which integrates both radios and scans Wi-Fi networks (channels 1--13) and Bluetooth Low Energy (BLE) beacons in dedicated switchable modes--- with a \textbf{real-time server} built on Node.js, Express and WebSocket that aggregates, analyses, persists and visualises the captured data on a web dashboard. EILOR provides four operating modes, an adjacency-aware channel-interference score, evil-twin access-point detection, OUI-based vendor identification, randomized-MAC awareness, and a token-based security layer. End-to-end validation confirmed data ingestion, channel analysis, restart-safe persistence, access control and captive-portal guest registration with offline buffering.
